
\section{4. Method}

\subsection{4.1 Overview}

A3 compresses a pretrained transformer in \textbf{five phases}:

1. \textbf{Calibration} ( 4.2): collect activation statistics via streaming hooks.
2. \textbf{Weighted SVD} ( 4.3): factorize each weight matrix in weighted space.
3. \textbf{Validation} ( 4.4): compute Frobenius reconstruction error.
4. \textbf{Adaptive escalation} ( 4.4.2): increase rank if error exceeds $\tau$.
5. \textbf{LayerNorm recalibration} ( 4.5): absorb distributional shift.

\textbf{Figure 4.1} shows the pipeline. Phase 1 runs once on the whole model; 
phases 2--4 run per layer; phase 5 runs once at the end.

\begin{lstlisting}
    +--------------------------------------------------+
    |  Phase 1: Calibration (streaming)                |
    |  Input: unlabeled text                           |
    |  Output: {F^( ), C_X^( )} per layer              |
    +----------------------+---------------------------+
                           |
                            
    +--------------------------------------------------+
    |  For each layer  :                               |
    |    Phase 2: Weighted SVD at rank r               |
    |    Phase 3: Frobenius validation                 |
    |    Phase 4: If err > $\tau$: escalate rank, retry     |
    |    If err > $\tau$ after max retries: fallback        |
    +----------------------+---------------------------+
                           |
                            
    +--------------------------------------------------+
    |  Phase 5: LayerNorm recalibration (global)       |
    |  Rescale $\gamma$ to match target $\sigma$                     |
    +--------------------------------------------------+
\end{lstlisting}

\subsection{4.2 Phase 1: Streaming Calibration}

\subsubsection{4.2.1 Problem}

We want to compute, for each linear layer $\ell$:

\[
F^{(\ell)} = \frac{1}{N} \sum_{t=1}^{N} \left(x_t^{(\ell)}\right)^2, \qquad
C_X^{(\ell)} = \frac{1}{N} \sum_{t=1}^{N} x_t^{(\ell)} \left(x_t^{(\ell)}\right)^\top
\tag{4.1}
\]

where $x_t^{(\ell)} \in \mathbb{R}^{d_\text{in}^{(\ell)}}$ is the input to 
layer $\ell$ at token $t$.

Naively, storing all $N$ activations per layer costs $\mathcal{O}(N \sum_\ell d_\text{in}^{(\ell)})$, 
which for $N = 10^4$ and $d_\text{in} = 768$ over 12 layers is $\sim 9 \times 10^7$ 
values, or 180 MB at FP16. Streaming avoids this.

\subsubsection{4.2.2 Streaming Accumulation}

We register forward hooks on every linear-like module. For each forward pass 
with input $x \in \mathbb{R}^{B \times L \times d}$, the hook:

1. Reshapes $x$ to $(BL) \times d$.
2. Filters padding tokens: keeps only rows where $\|x_i\|_1 > 10^{-5}$.
3. Accumulates $C_\text{batch} = x_\text{valid}^\top x_\text{valid} \in \mathbb{R}^{d \times d}$ and $f_\text{batch} = \sum_i x_{\text{valid}, i}^2 \in \mathbb{R}^d$.
4. Adds these to CPU-side running buffers.

After all calibration passes:

\[
F^{(\ell)} = \frac{\sum_k f^{(k)}}{\sum_k N^{(k)}}, \qquad
C_X^{(\ell)} = \frac{\sum_k C^{(k)}}{\sum_k N^{(k)}}
\tag{4.2}
\]

\textbf{Memory.} $\mathcal{O}(d^2)$ per layer, or $\mathcal{O}(N_l d^2)$ total. 
For $d = 768$ and $N_l = 12$, this is $\sim 8.5$M values, or 17 MB at FP32.

\subsubsection{4.2.3 Calibration Sources}

We support three sources, in priority order:

1. \textbf{Explicit \texttt{calibration\_loader}.} A \texttt{DataLoader} of input batches.
2. \textbf{\texttt{calibration\_texts}.} A list of strings; tokenized on-the-fly.
3. \textbf{Default.} Random inputs of the appropriate shape (weak calibration).

If none are provided, A3 falls back to unweighted SVD (see  4.3.4).

\subsection{4.3 Phase 2: Weighted SVD}

\subsubsection{4.3.1 The Weighting Idea}

Given weight $W \in \mathbb{R}^{d_\text{out} \times d_\text{in}}$, standard 
SVD minimizes $\|W - W_r\|_F^2$. We instead minimize an \textbf{activation-weighted} 
error:

\[
\mathcal{L}_\text{weighted}(W_r) = \mathbb{E}_x\!\left[\|W x - W_r x\|_2^2\right]
\tag{4.3}
\]

If we assume $x$ has approximately diagonal covariance 
$\mathrm{diag}(F_1, \dots, F_{d_\text{in}})$, then:

\[
\mathcal{L}_\text{weighted}(W_r) = \sum_{i=1}^{d_\text{in}} F_i \|(W - W_r)_{:, i}\|_2^2 = \|(W - W_r) \cdot \mathrm{diag}(\sqrt{F})\|_F^2
\tag{4.4}
\]

So the activation-weighted problem is equivalent to an \textbf{unweighted} problem 
on the transformed matrix:

\[
W_\text{weighted} = W \cdot \mathrm{diag}(\sqrt{F + \epsilon})
\tag{4.5}
\]

where $\epsilon = 10^{-6}$ prevents division by zero.

\subsubsection{4.3.2 The Algorithm}

1. Compute $W_\text{weighted}$ using (4.5).
2. Compute the truncated SVD:
\[
   W_\text{weighted} \approx U_r \Sigma_r V_r^\top
   \tag{4.6}
\]
3. Unweight the right factor to recover the original space:
\[
   V_r^\text{unweighted} = V_r^\top \cdot \mathrm{diag}\!\left(\frac{1}{\sqrt{F + \epsilon}}\right)
   \tag{4.7}
\]
4. Set $U = U_r \Sigma_r^{1/2}$ and $V = \Sigma_r^{1/2} V_r^\text{unweighted}$.

The factorized layer computes:

\[
y = (x V^\top) U^\top + b
\tag{4.8}
\]

\subsubsection{4.3.3 Scale Normalization}

To prevent extreme weighting values, we optionally normalize $F$:

\[
F_\text{norm} = \frac{F}{\|F\|_2 / \sqrt{d_\text{in}}}
\tag{4.9}
\]

and clamp the square root to $[10^{-3}, 10^3]$. This is controlled by the 
\texttt{normalize\_scale} flag.

\subsubsection{4.3.4 Fallback to Unweighted SVD}

If no calibration data is available, we set $F = \mathbf{1}$, giving 
standard SVD. This is a safe default but sacrifices the activation-aware 
benefit.

\subsection{4.4 Phase 3 \& 4: Validation and Adaptive Escalation}

\subsubsection{4.4.1 Frobenius Error Validation}

After factorization, we compute:

\[
\text{err}(W, U, V) = \frac{\|W - U V\|_F}{\|W\|_F + 10^{-8}}
\tag{4.10}
\]

If $\text{err} \leq \tau$, the layer is replaced. Otherwise:

\subsubsection{4.4.2 Adaptive Rank Escalation}

We retry with an increased rank:

\[
r_{k+1} = r_k + \max\!\left(8, \frac{\min(d_\text{in}, d_\text{out}) - r_k}{K - k}\right)
\tag{4.11}
\]

where $K$ is the maximum number of retries (default 3). The step size is 
chosen so the rank reaches $\min(d_\text{in}, d_\text{out})$ in exactly $K$ 
steps if no improvement is seen.

If even at maximum rank the error exceeds $\tau$, the layer **falls back to 
the original** (no compression), and we log the failure.

\subsubsection{4.4.3 Why Adaptive?}

Choosing rank upfront is a trade-off:

\begin{itemize}
\item Too high $\to$ poor compression
\end{itemize}
\begin{itemize}
\item Too low $\to$ excessive error
\end{itemize}

Our approach defers the decision to per-layer validation, which handles the 
natural variance in layer importance.

\subsection{4.5 Phase 5: LayerNorm Recalibration}

\subsubsection{4.5.1 Motivation}

After factorizing a layer, its output distribution changes. The LayerNorm 
that follows now sees a different input distribution, which may have 
non-unit variance even after normalization.

The LayerNorm computes:

\[
\text{LN}(x) = \gamma \odot \frac{x - \mu}{\sigma(x) + \epsilon} + \beta
\tag{4.12}
\]

If the input variance changed from $\sigma_\text{old}$ to $\sigma_\text{new}$, 
the output scaled by $\gamma / \sigma_\text{new}$ instead of $\gamma / \sigma_\text{old}$. 
We can compensate by adjusting $\gamma$:

\[
\gamma \leftarrow \gamma \cdot \frac{\sigma_\text{target}}{\sigma_\text{current}}
\tag{4.13}
\]

where $\sigma_\text{target}$ is the pre-compression standard deviation and 
$\sigma_\text{current}$ is the post-compression one.

\subsubsection{4.5.2 Procedure}

1. Before compression, record $\sigma_\text{target}^{(\ell)}$ for every 
   LayerNorm via a calibration pass.
2. After compression, run another calibration pass to compute $\sigma_\text{current}^{(\ell)}$.
3. For each LayerNorm, update $\gamma$ via (4.13).

This is applied only if \texttt{recalibrate\_ln = True} and calibration sources are 
available.

\subsection{4.6 Combined Algorithm}

\textbf{Algorithm 1:} A3 compression
\begin{lstlisting}
Input:  model M, rank r , threshold $\tau$, calibration data D
Output: compressed model M', statistics

 1: for   in linear_layers(M) do
 2:     F^( ), C_X^( ) $\leftarrow$ running_accumulate( , D)
 3: end for
 4: $\sigma$_target $\leftarrow$ record_layernorm_scales(M, D)
 5:
 6: for   in linear_layers(M) do
 7:     W $\leftarrow$ weight( )
 8:     r $\leftarrow$ r 
 9:     for attempt in 0..K do
10:         W_weighted $\leftarrow$ W   diag($\sqrt{}$(F^( ) + $\epsilon$))
11:         U_r,  _r, V_r $\leftarrow$ svd(W_weighted, r)
12:         V_r $\leftarrow$ unweight(V_r, F^( ))
13:         err $\leftarrow$  W - U_r V_r _F /  W _F
14:         if err $\le$ $\tau$: break
15:         r $\leftarrow$ escalate(r, K, attempt)
16:     end for
17:     if err > $\tau$: fallback( ) else: replace( , U_r, V_r)
18: end for
19:
20: if recalibrate_ln:
21:     $\sigma$_current $\leftarrow$ record_layernorm_scales(M', D)
22:     for ln in layernorms(M') do
23:         ln.$\gamma$ $\leftarrow$ ln.$\gamma$   $\sigma$_target[ln] / $\sigma$_current[ln]
24:     end for
25: end if
26: return M'
\end{lstlisting}

\subsection{4.7 Complexity Analysis}

\begin{table}[H]
\centering
\begin{tabular}{lll}
\toprule
\textbf{Phase} & \textbf{Time} & \textbf{Memory} \\
\midrule
Calibration & $\mathcal{O}(N \sum_\ell d_\ell^2)$ & $\mathcal{O}(N_l d^2)$ \\
Weighted SVD & $\mathcal{O}(\sum_\ell d_\text{out}^{(\ell)} d_\text{in}^{(\ell)} \min(d_\text{out}, d_\text{in}))$ & $\mathcal{O}(d_\text{max}^2)$ \\
Validation & $\mathcal{O}(\sum_\ell r d_\text{out} d_\text{in})$ & $\mathcal{O}(d_\text{out} d_\text{in})$ \\
LN recalibration & $\mathcal{O}(N)$ & $\mathcal{O}(N_l d)$ \\
\bottomrule
\end{tabular}
\end{table}

\textbf{Full compression time} on DistilBERT (66.9M params): ~{{COMPRESS\_TIME}} 
seconds on RTX 4070.

\subsection{4.8 Parameter Analysis}

For a layer with weight $W \in \mathbb{R}^{d_\text{out} \times d_\text{in}}$:

\begin{table}[H]
\centering
\begin{tabular}{ll}
\toprule
\textbf{Configuration} & \textbf{Parameters} \\
\midrule
Original & $d_\text{out} \cdot d_\text{in}$ \\
A3 at rank $r$ & $r(d_\text{out} + d_\text{in})$ \\
A3 at rank $r$ + LoRA($r_L$) & $r(d_\text{out} + d_\text{in}) + r_L(d_\text{out} + d_\text{in})$ \\
Reduction factor & $\frac{d_\text{out} d_\text{in}}{r(d_\text{out} + d_\text{in})}$ \\
\bottomrule
\end{tabular}
\end{table}

\textbf{Example.} For $d = 768$ and $r = 64$:

\begin{table}[H]
\centering
\begin{tabular}{ll}
\toprule
\textbf{Configuration} & \textbf{Params} \\
\midrule
Original & 589,824 \\
A3 rank 64 & 98,304 \\
A3 rank 64 + LoRA rank 16 & 122,880 \\
Reduction & **4.8$\times$** \\
\bottomrule
\end{tabular}
\end{table}

\subsection{4.9 Theoretical Notes}

\subsubsection{4.9.1 Optimality of Weighted SVD}

\textbf{Proposition.} The truncated SVD of $W_\text{weighted} = W \cdot \mathrm{diag}(\sqrt{F})$ 
minimizes $\|(W - W_r) \mathrm{diag}(\sqrt{F})\|_F$ over all rank-$r$ matrices 
$W_r$.

\textbf{Proof.} This follows directly from Eckart--Young applied to the matrix 
$W_\text{weighted}$. $\square$

\subsubsection{4.9.2 Relation to Standard SVD}

If $F = \mathbf{1}$ (uniform), then $W_\text{weighted} = W$ and the weighted 
SVD reduces to standard SVD.

\subsubsection{4.9.3 Error Bound}

For any $W$, $U_r$, $V_r$ from weighted SVD:

\[
\|W - U_r V_r\|_F \leq \left\|W - W^*_r\right\|_F + \|W^*_r - U_r V_r\|_F
\tag{4.14}
\]

where $W^*_r$ is the unweighted optimal rank-$r$ approximation. In practice, 
the weighted SVD often achieves lower $\|W - U_r V_r\|_F$ than unweighted 
because the weighting suppresses reconstruction error in low-variance 
directions.